\section{A terminating finite certificate}
\label{sec:finite}

For an orbit remaining after the listed channels have been removed,
the bound $D\le100$ is now a theorem, not a search parameter.
This section specifies the complete finite task and records its
accepted exact evaluation. The two enumerations were performed in
the proof stage; they were not rerun merely to prepare this article.
Appendix~\ref{app:certificate} prints the primary implementation
and identifies the preserved source and output files.

\subsection{The forward, positive-extremum algorithm}
For each integer $D\in[1,100]$ put
\[
 b_-=-3D-1,\qquad b_+=3D-1.
\]
Perform the following finite loops and exact orbit walks.
\begin{enumerate}
\item For every integer $u\in[b_-,b_+-D]$ and every
$s\in\{-3,-2,-1,0,1\}\cap[b_-,b_+]$, put $v=u+D$.
For every integer $p\in[-D,D]$, put $q=(s+1)D-p$.
Retain the seed precisely when
\begin{equation}
 |q|\le D,\qquad |(u+1)p|\le2D,\qquad
 |(v+1)q|\le2D.
 \label{eq:seedbounds}
\end{equation}
Set $a=v-su+s-p$ and the starting triple $P=(u,s,v)$.
The code intersects the corresponding integer $p$ intervals
before entering their innermost loop; this is exactly the same
finite set.
\item Keep this $a,D$ fixed. At the current state $(x,y,z)$,
test first whether a coordinate is outside $[b_-,b_+]$,
then whether $|z-x|>D$. If either test holds, record the
corresponding exit and stop this walk.
\item Otherwise append $x$ to the scalar word and replace the
state by $(y,z,yz+a-x)$. If it equals $P$, record the first
return and stop; otherwise repeat the preceding step.
\item For each returned word, check the recurrence at every
wraparound phase and absence of any smaller scalar period.
Insert both the word and its reverse, identifying words only
by cyclic rotation. Store the pair consisting of $a$ and the
lexicographically least rotation.
\item Compare each resulting word, at every cyclic phase, with
the explicit family formulas in Table~\ref{tab:classification}.
For this subtraction allow their degenerate parameter values;
classify by the actual word length. Output every unmatched word
with its parameter and first-return length.
\end{enumerate}

\begin{lemma}[Coverage and termination]
\label{lem:certificate}
The algorithm covers every orbit satisfying \eqref{eq:core},
including its original time orientation. Every orbit walk stops,
without any imposed bound on $a$ or on the number of iterates.
An exit excludes the seed from a periodic orbit having the
specified maximum difference and height bounds.
\end{lemma}
\begin{proof}
Every residual periodic word has a phase with $|d_0|=D$;
reverse if necessary to obtain $d_0=D$. Equation~\eqref{eq:key}
gives precisely the five choices of $s$, the relation defining
$q$, and the two neighboring bounds \eqref{eq:seedbounds}.
The height bounds place $u,v$ in the enumerated intervals.
The equation defining $p$ in \eqref{eq:extremal} uniquely
recovers $a$. Thus an extremal seed of this word occurs.
The reversed word is also a valid word for the same $a$ by
\eqref{eq:reversor}; explicitly restoring both reverses
recovers whichever orientation was used originally.

For fixed $a,D$, all nonexiting states belong to the finite
set $[b_-,b_+]^3\cap\Z^3$. If a walk never exited, it
would repeat a state. Under an injective map the first repeat
must be the initial state: if $P_j=P_i$ with $j>i>0$,
apply the inverse $i$ times to obtain $P_{j-i}=P_0$, an
earlier return. This proves termination at a first return or
exit. A genuine bounded periodic word could violate neither
necessary bound, so an exit is a valid exclusion. The
existence of this finite-state argument is not an operational
period cutoff in the code.
\end{proof}

\subsection{Independent interval seeds and inverse traversal}
The independent implementation changes the seed order, treats
both extremum signs directly, walks the inverse map, and compares
complete triple-state sets. It neither imports nor executes the
primary implementation. Its complete mathematical algorithm is
as follows.

For $D=1,\ldots,100$, $\delta\in\{-D,D\}$,
$s=-3,-2,-1,0,1$, and $q=-D,\ldots,D$, set
$p=(s+1)\delta-q$ and discard $|p|>D$.
Start with the integer interval
\begin{equation}
 I=[b_-,b_+]\cap[b_--\delta,b_+-\delta].
 \label{eq:interval0}
\end{equation}
If $p\ne0$, intersect $I$ with
\begin{equation}
 [-1-\lfloor2D/|p|\rfloor,-1+\lfloor2D/|p|\rfloor];
 \label{eq:intervalp}
\end{equation}
if $q\ne0$, intersect it with
\begin{equation}
 [-1-\delta-\lfloor2D/|q|\rfloor,
  -1-\delta+\lfloor2D/|q|\rfloor].
 \label{eq:intervalq}
\end{equation}
Zero $p$ or $q$ imposes no corresponding interval constraint.
For each integer $u\in I$, put
\begin{equation}
 v=u+\delta,\qquad a=q-sv+u+s,
 \qquad P=(u,s,v).
 \label{eq:independentseed}
\end{equation}
Check the reconstruction of both $p$ and $q$ and their bounds.
All five centers already lie in $[b_-,b_+]$ for $D\ge1$.
These interval intersections are exactly the integer solutions
of the same neighboring inequalities, now with both signs.

At each inverse-time state, test height first, then the difference
bound, and then membership in a visited set. A repeated state
must equal $P$; a noninitial first repeat raises an explicit
error. Otherwise insert the state and apply
$(x,y,z)\mapsto(xy+a-z,x,y)$, checking that a forward step
recovers the previous state. The proof of
Lemma~\ref{lem:certificate}, applied to the inverse bijection,
proves termination here too.

Upon return, choose the lexicographically least visited triple
and reconstruct its native \emph{forward} word. Check every
forward state belongs to the inverse visited set, that the first
repeat is the chosen triple, that all visited states are used,
and that the word has as many distinct triples as entries.
Represent this cycle by $(a,\min\{\text{visited triples}\})$,
while retaining the complete triple set. For fixed $a$ that set
determines the forward transitions uniquely, so this representation
does not quotient by reversal.

For independent family matching, generate candidate words rather
than use the primary rotation predicates. Generate $F_3$ with
$t=(a+4)/2$ when $a$ is even. For every value $t$ occurring
in the reconstructed word, generate $F_4$ and, at their required
parameters, $F_5,F_6,F_8$. At $a=0$ also generate $F_{12}$
with $m=\max_i|x_i|\ge1$. For a word of length at most
two, generate its alternating candidate. Verify candidate
recurrences and compare their complete triple-state sets.
Assign the first matching label in the order
$F_2,F_3,F_4,F_5,F_6,F_8,F_{12}$; this resolves the
degenerate family overlaps. An unmatched word is an exception.
These candidate rules cover every family match: its parameter
is either fixed by $a$, appears in the word, or equals its
maximum absolute coordinate in $F_{12}$.

Only after discovery is complete does the independent algorithm
parse the primary cycle output. It validates every recorded
integer word, recurrence, first-return length and family label,
rejects duplicate keys, compares the complete key sets in both
directions, and compares every full triple-state set. Reading
the input bytes earlier to hash them does not influence discovery.
Both implementations use exact unbounded integers.

\subsection{The complete finite outcome}
The preserved executions give Table~\ref{tab:certificate}.
Both implementations produce exactly $25\,851$ oriented cycles,
with the identical family partition
\begin{equation}
 (F_3,F_4,F_5,F_6,F_8,F_{12},\text{unmatched})
 =(200,25350,99,100,50,50,2).
 \label{eq:familypartition}
\end{equation}
These counts refer only to the finite certificate, not to the
infinite global families. There are $25\,753$ self-reversing
cycles and $49$ pairs of distinct reversed cycles in this output;
the latter all lie in $F_5$.
\begin{table}[htbp]
\centering\small
\caption{Complete historical seed dispositions in the proved
residual core. The extremum sign and traversal direction are
distinct choices. Every row is an exact partition; these are
accepted proof-stage executions, not new manuscript experiments.}
\label{tab:certificate}
\begin{tabular}{@{}lrrrr@{}}
\toprule
Traversal and extremum & Seeds & Returns & Difference exits & Height exits\\
\midrule
Forward, positive & 74,866 & 25,907 & 30,335 & 18,624\\
Inverse, negative & 74,866 & 25,907 & 30,335 & 18,624\\
Inverse, positive & 74,866 & 25,907 & 30,344 & 18,615\\
Inverse, both signs & 149,732 & 51,814 & 60,679 & 37,239\\
\bottomrule
\end{tabular}
\end{table}


The two unmatched words are
\begin{align*}
 a=-13 &: \quad (-4,-2,-3,-3,-2),\\
 a=-2 &: \quad (-2,-1,0,0,-1,-2,0,-1,0).
\end{align*}
Their first-return lengths are five and nine. The independently
reconstructed output agrees with every one of the primary
$25\,851$ records, not only its aggregate counts. The nine-seed
exchange between inverse exit types in
Table~\ref{tab:certificate} is not a mismatch: opposite
directions can first cross different rejecting boundaries.
The longest observed walks were $12$ forward and $18$ backward;
neither number was an input. Likewise the observed seed forcing
range $[-503,497]$ was an output, not a prescribed parameter
interval.

Corollary~\ref{cor:core}, Lemma~\ref{lem:certificate}, and
this complete exact residual outcome prove exhaustion of the
table. The remaining sections verify its existence, least-period,
disjointness and level-count assertions without relying on a
finite sample of its unbounded parameters.
